\section{Unrestricted local matching and its exact limitation}\label{sec:flexibility}
Finite local cohomological conditions need not obstruct a global character when its ramification is unrestricted. Neukirch's criterion, in its authenticated original form and the formulation recorded by Neftin, applies to both modules relevant to the two adjoint stages \cite{neukirch1973,neftin2013}.

For a finite module $A$, let $A'=\operatorname{Hom}(A,\mu_7)$ and let $G'$ be the finite group of the field cut out by that dual. The cited sufficient criterion is $H^1(G',A')=0$. For $A=k$ with trivial action, $A'\simeq k(\chi)$ and $G'=C_6$. Averaging kills its positive cohomology because six is invertible modulo seven. For $A=k(\chi)$, the source action cancels the cyclotomic target action, so the dual is trivial and $G'=1$. Both satisfy the criterion.

Consequently the unrestricted maps
\[
H^1(K,k)\longrightarrow\prod_{v\in S}H^1(K_v,k),\qquad
H^1(K,k(\chi))\longrightarrow\prod_{v\in S}H^1(K_v,k(\chi))
\]
are surjective for every finite set $S$ under the cited localization theorem. The precise source formulation and its fixed additional finite set refinement are recorded in the literature audit. This statement does not keep the character unramified at every other place.

If $p_0$ is one global bare primitive and $p_v$ are desired finite local primitives, their differences $p_v-p_0|_v$ are ordinary characters. For trivial coefficients there are no nonzero coboundaries. Surjectivity supplies one global character matching every prescribed finite difference exactly. Then $p=p_0+z$ meets those local equations, but may ramify at new primes.

For the upper equation, once one global $q_0$ exists, prescribed local differences are $k(\chi)$ cocycles. Their cohomology classes can likewise be matched at finitely many places. Local coboundaries correspond to strict upper conjugation. Thus conditions invariant under such conjugation cannot alone create an unrestricted finite localization obstruction when each is locally attainable.

\subsection{The local toric hyperplane}
At a beta support prime, use a pro seven tame generator $T$ and arithmetic Frobenius $F$. Put
\[
c=\beta(T)\ne0,\qquad b=\beta(F),\qquad h=(Q-1)/7\pmod7.
\]
The local first digit normal form is
\[
M=c\begin{pmatrix}\theta&-\theta^2\\1&-\theta\end{pmatrix},\qquad
A=\begin{pmatrix}\theta b&\theta h-\theta^2b\\b&h-\theta b\end{pmatrix}.
\]
Relative to the chosen reference,
\[
p(T)=\theta c,\qquad p(F)=\theta b-h.
\]
Eliminating $\theta$ gives the invariant affine equation
\begin{equation}\label{eq:toricp}
cp(F)-bp(T)=-hc.
\end{equation}
Conversely any pair satisfying this equation determines $\theta=p(T)/c$ and reconstructs the displayed first order packet. The packet converse does not construct a global curve or establish every higher integral relation.

Changing $F$ to $T^mF$ adds $mc$ to $b$ and $mp(T)$ to $p(F)$, so the left side of \eqref{eq:toricp} is unchanged. A consistent Frobenius lift change cannot remove the affine right side.

The apparent freedom from unrestricted localization is exactly why a justified ramification set is needed. Any useful obstruction must retain restrictions that the arbitrary global correction is not allowed to evade. These include the actual bad prime set, the one integral frame, and the next precision tame relation. Block 79 supplied such a set for the diagonal first cochain. It did not supply it for every upper correction.
